\needspace{12\baselineskip}
\begin{center}
\pdfbookmark[1]{8. Zu Art. 114 und 177.}{zusatz-008}
[8. Zu Art. 114 und 177.]
\end{center}
\begin{center}
\{Druckfehler in Dr. \textsc{Gauss}' Theoria motus corporum coelestium etc. Hamburgi 1809, vom Senator Bar. \textsc{Oriani}.\}
\end{center}
\begin{center}
\begin{tabular}{c}\hline
\small Monatliche Correspondenz Band XXI, S. 282--283, 1810 März.\\\hline
\end{tabular}
\end{center}
\{.... Ut aequatio 7 [art. 114] fiat identica, debent coëfficientes ipsorum \(\delta\delta'\delta''\), \(D\delta'\delta''\), \(DD'\delta''\) etc. esse \(=0\); sed post debitas reductiones coefficiens ipsius \(\delta\delta'\delta''\) est
\SourceLeaf{308}
\[
(0.1.2)=\left\{\begin{aligned}
&\operatorname{tang}\beta'\operatorname{tang}\beta\,\sin(L''-\alpha'')\sin(L'-L)\\
&+\operatorname{tang}\beta'\operatorname{tang}B\,\sin(L''-\alpha'')\sin(L'-\alpha)\\
&+\operatorname{tang}\beta\operatorname{tang}\beta''\sin(L'-\alpha')\sin(L-L'')\\
&+\operatorname{tang}B\operatorname{tang}\beta''\sin(L'-\alpha')\sin(\alpha-L'')\\
&+\operatorname{tang}\beta''\operatorname{tang}\beta'\sin(L-\alpha)\sin(L''-L')\\
&+\operatorname{tang}\beta''\operatorname{tang}B'\sin(L-\alpha)\sin(L''-\alpha')\\
&+\operatorname{tang}B'\operatorname{tang}\beta\,\sin(L''-\alpha'')\sin(\alpha'-L)\\
&+\operatorname{tang}B'\operatorname{tang}B\,\sin(L''-\alpha'')\sin(\alpha'-\alpha)\\
&+\operatorname{tang}\beta\operatorname{tang}B''\sin(L'-\alpha')\sin(L-\alpha'')\\
&+\operatorname{tang}B\operatorname{tang}B''\sin(L'-\alpha')\sin(\alpha-\alpha'')\\
&+\operatorname{tang}B''\operatorname{tang}\beta'\sin(L-\alpha)\sin(\alpha''-L')\\
&+\operatorname{tang}B''\operatorname{tang}B'\sin(L-\alpha)\sin(\alpha''-\alpha').
\end{aligned}\right.
\]
Hinc si habeatur tantummodo
\[
\left.\begin{aligned}
&\operatorname{tang}\beta'\operatorname{tang}\beta\,\sin(L''-\alpha'')\sin(L'-L)\\
&+\operatorname{tang}\beta\operatorname{tang}\beta''\sin(L'-\alpha')\sin(L-L'')\\
&+\operatorname{tang}\beta''\operatorname{tang}\beta'\sin(L-\alpha)\sin(L''-L')
\end{aligned}\right\}=0,
\]
idem coefficiens non fit \(=0\), sed requiritur praeterea ut sit
\[
B=0,\qquad B'=0,\qquad B''=0.\}
\]
\begin{center}\rule{.25\textwidth}{.35pt}\end{center}

\{Elegans theorema, quod tribuitur [art. 117] illustr. \textsc{Laplace}, revera a \textsc{Leonardo Eulero} primum inventum est. Et enim in Comment. Acad. Petropol. Tom. XVI\HistoricalNoteEight{Handschriftliche Bemerkung von \textsc{Gauss}:] \textls[150]{Dies} Theorem findet sich a. a. O. nicht, wohl aber p. 101 folgendes: \(\int_0^1\BareRoot{\left(\log\frac1x\right)}\mathrm{d}x=\tfrac12\BareRoot{\pi}.\) Schreibt man hier \(x=e^{-tt}\), so wird \(\int_0^\infty2tte^{-tt}\mathrm{d}t=\tfrac12\BareRoot{\pi}.\) Es ist aber \(2tte^{-tt}\mathrm{d}t=-\mathrm{d}(te^{-tt})+e^{-tt}\mathrm{d}t\) und \(te^{-tt}=0\) sowohl für \(t=0\) als für \(t=\infty\), also \(\int_0^\infty e^{-tt}\mathrm{d}t=\tfrac12\BareRoot{\pi}.\)} \textsc{Eulerus} ostendit, integrale \(-\int\frac{\mathrm{d}x}{\BareRoot{\left(\log\frac1x\right)}}\) sumtum ab \(x=1\) ad \(x=0\) esse \(=\BareRoot{\pi}\), existente \(\pi\) semicircumferentia circuli, radio \(=1\) descripti. Iamvero ponendo \(x=e^{-tt}\) habetur
\[
\frac{-\mathrm{d}x}{\BareRoot{\left(\log\frac1x\right)}}=2e^{-tt}\mathrm{d}t.
\]

Ideoque integrale \(\int e^{-tt}\mathrm{d}t\) a \(t=0\) ad \(t=\infty\) erit \(=\tfrac12\BareRoot{\pi}\), et propterea idem integrale a \(t=-\infty\) ad \(t=+\infty\) fiet \(=\BareRoot{\pi}.\)\}
\begin{center}\rule{.25\textwidth}{.35pt}\end{center}

\SourceLeaf{309}
\begin{center}
\begin{tabular}{c}\hline
\small Monatliche Correspondenz. Band XXI. S. 280. 1810 März.\\\hline
\end{tabular}
\end{center}
Für die Notirung der Druckfehler in meiner Theoria bin ich Herrn \textsc{Oriani} sehr verbunden. Er hat ganz Recht, dass ich pag. 129 [Art. 114] hinzuzufügen vergessen habe, dass \(B,B',B''=0\) vorausgesetzt werden müssen, wenn die Bedingungsgleichung, bei welcher die Gleichung 7 unbrauchbar ist, die dort angegebene Gestalt haben soll. Es ist übrigens klar, dass wenn auch nicht \(B,B',B''=0\) sind, doch die Gleichung 7 unbrauchbar sein kann, wenn nemlich der 12-gliedrige Ausdruck, welchen \textsc{Oriani} entwickelt hat, zufällig \(=0\) oder sehr klein wird.

Dass \textsc{Euler} schon das Theorem gefunden hat, woraus der schöne von mir \textsc{Laplace} beigelegte Lehrsatz sehr leicht abgeleitet werden kann, fiel mir selbst schon früher ein, als aber die Stelle pag. 212 [Art. 177] schon abgedruckt war; ich wollte es aber nicht unter die Errata setzen, weil \textsc{Laplace} wenigstens das obige Theorem doch erst in der dort gebrauchten Form aufgestellt hat. ....
\begin{center}\rule{.18\textwidth}{.35pt}\end{center}
